\clearpage
\unitheader{D}{Part 4: Initialization, feature learning, and transfer}{}

So far, we have relied on a scaling-law based approach to finding good hyperparameters at large scale. However, even a theory-guided approach of this kind may seem unsatisfying. Why does the optimal hyperparameter have to change at all? If we design the architecture and optimizer right, perhaps we could tune a hyperparameter once, and use that same optimal hyperparameter everywhere. This is the promise of various hyperparameter transfer schemes, and we will go in-depth into a popular kind of transfer scheme known as $\mu P$.

\subunitheader{The goal and the $abc$ parameterization}
We want to tune hyperparameters on a small model and reuse them as we scale width
and depth. It will turn out that we can patch up existing architectures and optimizers to have reasonable hyperparameter transfer by adjusting some small details in how we parametrize initialization, output scale, and learning rates.

Following
\href{https://arxiv.org/abs/2407.05872}{Everett et al.}, consider a one linear
layer with fan-in $n$, $p$ outputs, and input feature $x\in\mathbb R^n$.
We adjust the standard form of this layer with three exponents $(a,b,c)$ which modify the forward multiplier, initialization scale,
and learning rate. Let $\sigma$ and base LR $\eta$ be scale independent constants (i.e. our hyperparameters that are stable with scale),
\[
 z=n^{-a}Wx,\qquad W\in\mathbb R^{p\times n},\qquad
 (W_0)_{ij}\overset{\rm iid}{\sim}\mathcal N(0,\sigma^2n^{-2b}),\qquad
 \eta_W=\eta n^{-c}.
\]

\subunitheader{Two principles: stability and feature learning}
We compare per-coordinate scales using RMS norms:
\[
 \|v\|_{\rm RMS}^2=\frac1n\sum_{j=1}^nv_j^2,\qquad
 \|A\|_{\rm RMS}^2=\frac1{pn}\sum_{i,j}A_{ij}^2
 \quad(A\in\mathbb R^{p\times n}).
\]
For output vectors in $\mathbb R^p$, replace $n$ by $p$ in the vector RMS definition.
Take width $n\to\infty$ at a fixed training step $t$. For a fixed probe input
$x$ with $\|x\|_{\rm RMS}=\Theta(1)$, let $z_t=n^{-a}W_tx$.
Here $O(1)$ means bounded with width and $\Theta(1)$ means bounded and nonvanishing.
\emph{Stability} requires usable initial hidden features and bounded changes:
\[
 \|z_0\|_{\rm RMS}=\Theta(1),\qquad
 \|z_t-z_0\|_{\rm RMS}=O(1).
\]
Initial readout logits may instead vanish, so their initial scale need only be $O(1)$. The $\Theta(1)$ initialization requirement above applies to hidden features.
\emph{Feature learning} requires a nonvanishing change at some fixed step:
\[
 \|z_t-z_0\|_{\rm RMS}
 =\|n^{-a}(W_t-W_0)x\|_{\rm RMS}=\Theta(1).
\]
If this change scales as $n^{-r}$, then $r<0$ is unstable, $r=0$ gives feature
learning, and $r>0$ gives a vanishing change.

\subunitheader{Deriving Kaiming initialization from stability}
We first use stability at initialization to derive Kaiming initialization.
Since $(W_0)_{ij}\overset{\rm iid}{\sim}\mathcal N(0,\sigma^2n^{-2b})$
and $W_0$ is independent of the input $x$, the cross terms vanish:
\[
 \mathbb E[z_{0,i}^2\mid x]
 =n^{-2a}\sum_{j=1}^n\mathbb E[(W_0)_{ij}^2]x_j^2
 =\sigma^2n^{-2(a+b)}\|x\|_2^2.
\]
For stability, we want all hidden features, including $x$ and $z$, to have
$\Theta(1)$ RMS norm at initialization. This motivates choosing $a+b=1/2$:
\[
 a+b=1/2\quad\Longrightarrow\quad
 \mathbb E[\|z_0\|_{\rm RMS}^2\mid x]
 =\sigma^2\|x\|_{\rm RMS}^2.
\]
Kaiming initialization for ReLU chooses $a=0$, $b=1/2$, and $\sigma^2=2$.
See \href{https://openaccess.thecvf.com/content_iccv_2015/html/He_Delving_Deep_into_ICCV_2015_paper.html}{He et al.\ (2015), Section 2.2}
for the derivation of the constant factor $2$.

Kaiming initialization does not address feature learning during training,
which motivates $\mu$P theory. To understand it, we introduce notation for
alignment between the weight update and the input feature.

\subunitheader{Alignment between the weight update and the input feature}
Let $\Delta W := W_{t}-W_{0}$. The training pair of interest is $(\Delta W,x)$. By definition, a weight update gives $\Delta z=n^{-a}\Delta W x$. For nonzero
$\Delta W$ and $x$, define
\[
 R(\Delta W,x)=\frac{\|\Delta W x\|_{\rm RMS}}
 {\|\Delta W\|_{\rm RMS}\|x\|_{\rm RMS}},\qquad
 \|\Delta z\|_{\rm RMS}
 =n^{-a}R(\Delta W,x)\|\Delta W\|_{\rm RMS}\|x\|_{\rm RMS}.
\]
The ratio $R$ represents the alignment between the parameter update and the inputs. Note how the alignment allows us to simplify the change in features in terms of their constituent parts (change in parameters and change in inputs). Now let us get some intuitions for the order of magnitude scale for the feature updates in two special cases.

Assume $\Delta W\in\mathbb R^{p\times n}$, where $n$ is the fan-in.

\paragraph{No alignment: $\Delta W$ has iid centered Gaussian entries and is independent of $x$, so $R(\Delta W,x)\approx\sqrt{n}$.}
Suppose the update entries $\Delta W_{ij}$ are independent
$\mathcal N(0,\tau^2)$ variables, independent of $x$.
For $j\ne r$, $\mathbb E[\Delta W_{ij}\Delta W_{ir}]=0$, so
\[
 \mathbb E[(\Delta W x)_i^2\mid x]
 =\sum_{j,r}x_jx_r\mathbb E[\Delta W_{ij}\Delta W_{ir}]
 =\tau^2\sum_jx_j^2 =n\tau^2\|x\|_{\rm RMS}^2.
\]
Thus the RMS scale of $\Delta W x$ is $\sqrt{n}\,\tau\|x\|_{\rm RMS}$,
while that of $\Delta W$ is $\tau$. Gaussian rotational symmetry gives
the exact RMS scale of the ratio:
\[
 \bigl(\mathbb E[R(\Delta W,x)^2\mid x]\bigr)^{1/2}=\sqrt{n}.
\]
The equality holds in RMS over random updates; an individual update need not
have exactly this ratio.

\paragraph{Full alignment: $\Delta W$ has rank one and $x$ lies along its right singular vector, so $R(\Delta W,x)=n$.}
A gradient built from its input has this form. For $t=1$, take one SGD step from $W_0$ on a single-example loss $\ell(z)$,
where $z=n^{-a}W_0x$ and $\delta=\partial\ell/\partial z$ is the loss gradient
with respect to the output. With learning rate $\eta_W$, the update is
\[
 \Delta W=-\eta_W n^{-a}\delta x^\top,\qquad
 \Delta W x=-\eta_W n^{-a}\delta\|x\|_2^2.
\]
For a nonzero update, $\Delta W$ has rank one and $x/\|x\|_2$ is its
right singular vector with nonzero singular value. Each row is proportional
to $x$, so we have
\[
 \|\Delta W\|_{\rm RMS}
 =\eta_W n^{-a}\|\delta\|_{\rm RMS}\|x\|_{\rm RMS},\qquad
 R(\Delta W,x)=\frac{\|x\|_2^2}{\|x\|_{\rm RMS}^2}=n.
\]
This attains the Cauchy--Schwarz bound $R(\Delta W,x)\le n$.

We see that $\sqrt{n}$ and $n$ are both reasonable scalings.
We will see later which of these is empirically a better fit.


\needspace{20\baselineskip}
\subunitheader{Alignment between the readout head and its input feature change}
Let $V_0\in\mathbb R^{n\times q}$ be the initial readout matrix in the
transposed convention, where the number of classes $q$ does not scale with $n$.
For a last-layer feature change $\Delta x\in\mathbb R^n$, the initial
readout contributes $\Delta z_{\rm feature}=n^{-a}V_0^\top\Delta x$.
For nonzero $V_0$ and $\Delta x$, define
\[
 S(V_0,\Delta x)=\frac{\|V_0^\top\Delta x\|_{\rm RMS}}
 {\|V_0\|_{\rm RMS}\|\Delta x\|_{\rm RMS}},\qquad
 \|\Delta z_{\rm feature}\|_{\rm RMS}
 =n^{-a}S(V_0,\Delta x)\|V_0\|_{\rm RMS}\|\Delta x\|_{\rm RMS}.
\]
There are two reference cases:
\begin{itemize}
\item \textbf{No alignment:} iid centered Gaussian $V_0$ independent of $\Delta x$
gives $\bigl(\mathbb E[S(V_0,\Delta x)^2\mid\Delta x]\bigr)^{1/2}=\sqrt n$.
\item \textbf{Full alignment:} the columns of $V_0$ are proportional to $\Delta x$,
giving $S(V_0,\Delta x)=n$.
\end{itemize}

\subunitheader{Deriving $\mu$P by layer type}
For this derivation, we assume that both full-alignment conditions hold simultaneously:
\begin{enumerate}
\item Full alignment between the weight update and the input feature.
\item Full alignment between the initial readout head and the input feature change.
\end{enumerate}
Problem~4.0 asks you to derive the scaling rules for all four combinations of
these alignment assumptions.

\needspace{10\baselineskip}
\paragraph{Hidden and readout share the weight-update constraint.}
Maximal-update parameterization ($\mu$P) keeps updates nonvanishing and signals bounded.
Write an Adam update as $\Delta W=-\eta n^{-c}U$, where
$U=\widehat m/(\sqrt{\widehat v}+\epsilon)$ is the normalized direction.
For Adam, we can assume $\|U\|_{\rm RMS}=\Theta(1)$.\footnote{Adam can be viewed as signSGD with momentum, using $\operatorname{sign}(\widehat m)$, together with variance adaptation; see \href{https://arxiv.org/pdf/1705.07774}{Balles and Hennig (2018), \emph{Dissecting Adam}}.}
Define $\alpha$ by $R(U,x)=\Theta(n^\alpha)$. For an order-one input,
\[
 \|n^{-a}\Delta W x\|_{\rm RMS}
 =\eta n^{-a-c}R(U,x)\|U\|_{\rm RMS}\|x\|_{\rm RMS}
 =\Theta(\eta n^{\alpha-a-c}).
\]
For feature learning, we want $\|n^{-a}\Delta W x\|_{\rm RMS}$ to stay
at a constant, nonzero scale as $n\to\infty$. Under full alignment
($\alpha=1$), this gives our first equation:
\[
 a+c=1.
\]

\paragraph{The feature-change constraint depends on the layer's shape.}
After one update from $(W_0,x_0)$, the changes $\Delta W$ and $\Delta x$ give
\[
 \Delta z=n^{-a}\bigl(\underbrace{\Delta W x_0}_{\text{weight update}}
 +\underbrace{W_0\Delta x}_{\text{feature change}}
 +\underbrace{\Delta W\Delta x}_{\text{interaction}}\bigr).
\]
We have already shown that the direct weight-update term satisfies
$\|n^{-a}\Delta W x_0\|_{\rm RMS}=\Theta(1)$.
For $\|\Delta x\|_{\rm RMS}=\Theta(1)$, the interaction obeys
\[
 \|n^{-a}\Delta W\Delta x\|_{\rm RMS}
 \le n^{1-a}\|\Delta W\|_{\rm RMS}\|\Delta x\|_{\rm RMS}
 =O(\eta n^{1-a-c})=O(\eta).
\]
It remains to control $W_0\Delta x$. Assume $W_0$ has $p$ output channels.

To express this term's dependence on $n$, we use the Gaussian matrix norm bound.
Let $\|A\|_{\rm op}=\sup_{u\ne0}\|Au\|_2/\|u\|_2$ be the operator norm.
The
\href{https://www.math.uci.edu/~rvershyn/papers/HDP-book/HDP-1.pdf}{Gaussian matrix norm bound}
 gives $\|W_0\|_{\rm op}=O(n^{-b}(\sqrt p+\sqrt n))$ with high probability.
Converting to RMS norms bounds every $\Delta x$ with RMS $\Theta(1)$, including
changes that depend on $W_0$:
\[
 \|n^{-a}W_0\Delta x\|_{\rm RMS}
 \le n^{-a}\sqrt{\frac np}\,\|W_0\|_{\rm op}\|\Delta x\|_{\rm RMS}
 =O\!\left(n^{-a-b}\left(\sqrt n+\frac n{\sqrt p}\right)\right).
\]

\paragraph{1. Hidden matrices: both dimensions scale together, $p\propto n$.}
For $p\propto n$, the feature-change bound is $O(n^{1/2-a-b})$.
The initialization requirement $a+b=1/2$ already makes this $O(1)$, even when
$\Delta x$ depends on $W_0$. No stronger initialization scaling is needed.

\paragraph{2. Readout: full readout--feature-change alignment gives $a+b=1$.}
For the readout, $p=q$ is fixed. Write $V_0=W_0^\top\in\mathbb R^{n\times q}$
for the initial readout in the transposed convention. Full alignment with
the learned feature change means $S(V_0,\Delta x)=\Theta(n)$.
With $\|V_0\|_{\rm RMS}=\Theta(n^{-b})$ and
$\|\Delta x\|_{\rm RMS}=\Theta(1)$, this gives
\[
 \|n^{-a}V_0^\top\Delta x\|_{\rm RMS}
 =n^{-a}S(V_0,\Delta x)\|V_0\|_{\rm RMS}\|\Delta x\|_{\rm RMS}
 =\Theta(n^{1-a-b}).
\]
Stability requires $a+b\ge1$. Requiring a nonvanishing, order-one response
to the feature change gives $a+b=1$.
At initialization, the readout and its input feature are independent, so the
initial logits have RMS scale $\Theta(n^{1/2-a-b})=\Theta(n^{-1/2})$.

\paragraph{3. Input matrices: fan-in stays fixed.}
With fixed fan-in $d$ and unchanged external inputs, initialization and direct
update scales are $n^{-a-b}$ and $\eta n^{-a-c}$, assuming nonvanishing update
alignment. Keeping both order one gives $a+b=a+c=0$.

Our derivations show that different layers require different choices of
$(a,b,c)$ for stable, nonvanishing feature updates. We summarize one such
choice below. Here $m=n/n_0$, $n_0$ is the reference width, and $k$ is each
matrix's actual fan-in; the scalings match the baseline at $n_0$.
\[
 \begin{array}{c|ccc}
 &\text{input}&\text{hidden}&\text{readout}\\\hline
 (a,b,c)&(0,0,0)&(0,1/2,1)&(1,0,0)\\
 \text{forward multiplier}&1&1&1/m\\
 \text{initialization variance}&1/d&1/k&1/n_0\\
 \text{LR}&\eta&\eta/m&\eta
 \end{array}
\]
\begingroup
\setcounter{question}{-1}
\renewcommand{\aTwoProblemNumber}{4.\arabic{question}}
\needspace{20\baselineskip}
\begin{problem}{Derive scaling rules for update and feature-change alignment}
Consider all four combinations of two assumptions, each full or no alignment:
\begin{itemize}
\item Alignment of the weight-update direction $U$ with the pre-update input $x_0$.
\item Alignment of the initial readout head $V_0\in\mathbb R^{n\times q}$
with its input feature change $\Delta x=x_1-x_0$.
\end{itemize}
Here $V_0$ uses the transposed convention, so its response to the feature
change is $n^{-a}V_0^\top\Delta x$.
Write $\Delta W=-\eta n^{-c}U$, with fixed $\eta>0$ and
$\|U\|_{\rm RMS}=\Theta(1)$. For hidden and readout matrices, assume
$x_0$ and $\Delta x$ have order-one RMS, and write
\[
 R(U,x_0)=\Theta(n^\alpha),\qquad
 S(V_0,\Delta x)=\Theta(n^{\omega_{\rm move}}).
\]
For the interaction term, also assume $R(U,\Delta x)=O(n^\alpha)$.
Use $a=0$ for hidden and input matrices and $c=0$ for the readout.
\begin{parts}
\item For each combination, identify $\alpha$ and $\omega_{\rm move}$ and derive the
constraints on $(a,b,c)$. Require order-one hidden features at initialization,
order-one direct updates, and an order-one response $n^{-a}V_0^\top\Delta x$ from the initial readout.
Summarize $(a,b,c)$ for hidden, readout, and fixed-fan-in input matrices in a table.
\item Using $m=n/n_0$, give the forward multiplier, initialization variance,
and LR for each layer type, matching the preceding table at $n_0$.
How do the scaling rules change with each alignment assumption?
Check that the interaction term stays bounded and identify the combination
that recovers the preceding $\mu$P rules.
\end{parts}
\end{problem}

\clearpage
\subunitheader{4.1: Toy Example: A five-step Transformer stress test}

Test how $\mu$P scaling transfers hyperparameters across extreme widths and depths.

\paragraph{Training protocol.}
Train for five updates with constant LR from the first update, no warmup,
no weight decay, and no gradient clipping. Use unscaled rotary position
embeddings (RoPE). Use the same five training batches
and 64 fixed validation sequences across parameterizations; evaluate in
microbatches and tune LR using validation loss after update five.

\paragraph{Width settings.}
Let $m=n/n_0$ and let $k$ be each matrix's fan-in. Initialize weight matrices
with independent centered Gaussian entries using the variances below.
\begin{center}
\begin{tabular}{l@{\hspace{2em}}c@{\hspace{2em}}c}
\toprule
Setting & Kaiming & $\mu$P \\
\midrule
Depth & \multicolumn{2}{c}{2} \\
Attention-head dimension & \multicolumn{2}{c}{64} \\
Reference width $n_0$ & \multicolumn{2}{c}{512} \\
Weights, activations, optimizer states & \multicolumn{2}{c}{FP32} \\
Embedding variance & \multicolumn{2}{c}{1} \\
Hidden-matrix variance & \multicolumn{2}{c}{$1/k$} \\
Initial normalization gains & \multicolumn{2}{c}{1} \\
\midrule
Readout variance & $1/n$ & $1/n_0$ \\
Readout forward multiplier & $1$ & $1/m$ \\
Hidden-matrix LR & $\eta$ & $\eta/m$ \\
Embedding, readout, norm LR & $\eta$ & $\eta$ \\
\bottomrule
\end{tabular}
\end{center}

\paragraph{Depth settings.}
Fix width and reference width at 64, with one attention head and reference
depth two. Keep parameters, gradients, Adam states, and residual additions
in FP32; use BF16 autocast for branch computations.
Let $L$ count Transformer layers, $r=L/2$, and $c_L$ multiply each attention
and feedforward branch before residual addition. Apply these depth multipliers
on top of the $\mu$P width prescription:
\begin{center}
\begin{tabular}{lccc}
\toprule
 & Residual $c_L$ & Block LR multiplier & Block Adam $\epsilon$ \\
\midrule
Standard $\mu$P & $1$ & $1$ & $10^{-8}$ \\
\href{https://arxiv.org/abs/2310.02244}{Depth-$\mu$P} & $r^{-1/2}$ & $r^{-1/2}$ & $10^{-8}r^{-1/2}$ \\
\href{https://arxiv.org/abs/2505.01618}{CompleteP} & $r^{-1}$ & $1$ & $10^{-8}r^{-1}$ \\
\bottomrule
\end{tabular}
\end{center}
Block multipliers apply to hidden matrices and within-block normalization.
Embedding, readout, and final-normalization LR and stabilizer remain unchanged.

\needspace{12\baselineskip}
\begin{problem}{Test width and depth scaling over five updates (about 55 runs)}
\paragraph{Measure feature movement, update alignment, and readout--feature-change alignment.}
Let $h_t$ be a hidden feature vector of width $n$, evaluated on $N$ fixed
probes $x_1,\ldots,x_N$. For a matrix $A_t$ with fan-in $k$, let
$\Delta A_t=A_t-A_{t-1}$ be its actual update and $u_{t-1}$ its pre-update input.
Measure feature movement from initialization and update alignment by
\[
 M_t=\left[\frac1{Nn}\sum_{i=1}^N
 \|h_t(x_i)-h_0(x_i)\|_2^2\right]^{1/2},\qquad
 \alpha_{\rm upd}^{(t)}=\log_k R(\Delta A_t,u_{t-1}).
\]
Pool feature RMS over sequences, token positions, and coordinates; compute
weight RMS over entries. Inspect normalized features and the residual stream
separately.

For readout probes, take $h_t$ to be the final RMSNorm output before the
readout multiplier, and let $V_0\in\mathbb R^{n\times q}$ be the initial
readout weights in the transposed convention. Measure
\[
 \omega_{\rm move}^{(t)}=\log_n S(V_0,h_t-h_0).
\]
This probes the initial-readout/feature-change term in the derivation.
Use the same fixed probes and pool RMS as above; omit the ratio when
feature movement is zero.

\begin{parts}
\item \textbf{Transfer across widths (about 30 five-step runs).}
Predict how well Kaiming and standard $\mu$P will transfer their best base LR
across widths 640, 2,560, and 5,120. Choose an LR sweep for each width
and parameterization.
Fit each loss--LR curve with a suitable model. Plot fitted optimal base LR
and fitted minimum validation loss against width. Which parameterization
transfers its tuned base LR better? Does it also achieve lower loss?

\item \textbf{Probe width scaling (reuse part (a)'s configurations).}
At each width and parameterization's best sampled LR, track logit RMS and
the RMS and movement $M_t$ of the final RMSNorm output on fixed probes
through the five updates. Plot update alignment $\alpha_{\rm upd}^{(t)}$ for
each block's query, key, value, attention-output, and feedforward gate/up/down
matrices, and the final readout. Use each matrix's pre-update input and its
actual fan-in $k$. Also plot $\omega_{\rm move}^{(t)}$, comparing with the no-alignment ($1/2$) and
full-alignment ($1$) scales. Compare feature movement and alignment across
layers and widths. Which assumptions and predictions of the preceding derivation agree
with your measurements?

\needspace{7\baselineskip}
\item \textbf{Transfer across depths (about 25 five-step runs).}
Predict which depth prescription should transfer its base LR most reliably.
Compare standard $\mu$P, Depth-$\mu$P, and CompleteP at depths 2, 100, and
1,000, choosing an LR sweep for each depth and prescription.
Plot fitted optimal base LR and fitted minimum validation loss against depth.
Does better LR transfer imply lower loss after five updates?

\item \textbf{Probe depth scaling (reuse part (c)'s configurations).}
At each depth and prescription's best sampled LR, track residual-stream RMS
before normalization and branch-output RMS before multiplication by $c_L$.
Measure movement of the normalized residual-stream features at matched
relative positions, such as the midpoint and final layer.
Also track $\omega_{\rm move}^{(t)}$ through
all five updates and compare it across depths and prescriptions.
Can similar normalized activations hide a growing residual stream?
Relate these probes to the LR-transfer results in part (c), distinguishing
bounded signals from nonvanishing feature changes.
\end{parts}
\end{problem}

\clearpage
\subunitheader{4.2: Do the scaling rules predict longer training?}

Do the scaling ideas from Problem 4.1 help over longer training? Use the
default model and training settings at 153.6M tokens.

The course baseline uses the default architecture and initialization.
For $\mu$P, change only the four settings in the table below, using reference
width $n_0=512$. Let $n$ be hidden width, $m=n/n_0$, and $\eta$ the peak base LR.
Let $G$ be a shared matrix of standard Gaussian draws truncated at $\pm3$,
used for both embedding and readout initialization and shared across prescriptions.
\begin{center}
\begin{tabular}{lcc}
\toprule
Setting & Course baseline (default) & $\mu$P \\
\midrule
Embedding initialization & $G/n$ & $G/n_0$ \\
Readout initialization & $G/\sqrt n$ & $G/\sqrt{n_0}$ \\
Readout forward multiplier & $1$ & $1/m$ \\
Hidden-matrix peak LR & $\eta$ & $\eta/m$ \\
\bottomrule
\end{tabular}
\end{center}
Keep all other training settings fixed, including RMSNorm $\epsilon=10^{-5}$
and gradient clipping at 1. Normalization and clipping can affect widths
differently: the baseline embedding scale $1/n$ may be comparable to
$\sqrt\epsilon$. Use feature logs and \texttt{gradients.jsonl} to report
pre-clipping gradient norms, the embedding's fraction of squared gradient
norm, and the clipping coefficient. Account for these effects when interpreting
transfer. Reuse the supplied width-512 diagnostics; optional normalization or
clipping ablations must fit within your compute budget.


\begin{problem}{From five updates to a longer training budget (recommended: 50 new runs)}
\begin{parts}
\item \textbf{Transfer the reference LR (20 initial runs).}
Predict whether Problem 1's best sampled LR at 153.6M tokens transfers to
widths 128 and 256 under both prescriptions. Choose an LR sweep for each
configuration, including the transferred source LR.
Hold data and other settings fixed across new runs; compare with Problem 1's
supplied width-512 curve.

At the source LR, plot $\alpha_{\rm upd}^{(t)}$ throughout training for both
prescriptions and all three widths. Use fixed inputs and each matrix's fan-in
for every block's attention and feedforward matrices and the readout.
Also plot $\omega_{\rm move}^{(t)}$ from Problem 4.1. How do update alignment
and readout--feature-change alignment evolve beyond the first five updates?

\item \textbf{Compare transfer and tuned performance (no new runs).}
Fit each loss--LR curve. Plot fitted optimal base LR and fitted minimum validation
loss against width for both prescriptions. Also plot the measured loss at the
transferred LR. Does the prescription with better LR transfer also achieve lower
loss after tuning? How does the comparison differ from the five-step test?

\item \textbf{Predict a held-out width (about 12 new runs).}
Hold out width 1,024. For each prescription, fit a power law to the optimal base
LRs at widths 128, 256, and 512 from part (b), and predict its optimal LR at 1,024.
Train at the predicted LR and at the directly transferred source LR. Then run a
local LR sweep at 1,024. Compare both predicted LRs with the fitted local optimum and
report their measured loss gaps relative to the best sampled target LR. Does fitting the width dependence improve on
direct transfer?

\item \textbf{Transfer across depth (18 initial runs).}
Fix width 512 and compare depths 4, 8, and 16. Test standard $\mu$P,
Depth-$\mu$P, and CompleteP using the rules from Problem 4.1, now with reference
depth eight: for layer count $L$, use relative depth $r=L/8$.
At depths 4 and 16, compare the transferred source LR with an LR sweep of
your choice. Plot fitted optimal base LR and fitted minimum
validation loss against depth, alongside the measured losses at the transferred
LR. Does a depth correction improve transfer, tuned performance, or both?

\item \textbf{Identify the training regime (diagnostics from the same runs).}
At the directly transferred LR and the best sampled LR for each configuration,
track logit RMS and normalized hidden-feature movement on fixed inputs during
the first five updates and at the end of training, as in Problem 4.1. Include both the held-out width and the depth comparison.
Include $\omega_{\rm move}^{(t)}$ at the same
checkpoints. Which alignment assumptions are supported at different widths,
depths, and stages of training?
Can a configuration with larger early transients still achieve a better final loss?
Does a shift in optimal LR alone show that the model is outside the stable
$\mu$P regime? Use your diagnostics to explain what the experiments establish.
Use the LR--WD coupling from Problem 2 to discuss how holding WD fixed might
affect the comparison.
\item \textbf{Muon and $\mu$P (optional).}
\href{https://kellerjordan.github.io/posts/muon/}{Muon} is a momentum-based
optimizer that approximately orthogonalizes the update direction for each
hidden weight matrix before applying it. Explore how Muon interacts with
$\mu$P, following the same approach as in parts (a)--(c).
This exploration is not part of the required workload; attempt it only if you
have GPU budget remaining after the required experiments.

\end{parts}
\end{problem}

\endgroup
\setcounter{question}{4}

\clearpage
\begin{examplequestion}[difficulty=easy]
\phantomsection\label{ex:a2-copied-width}
\textbf{Widening a Kaiming model by copying.}
Use the five-step width-stress-test setting from Problem 4.1 on matching
training batches. Assume Adam's numerical stabilizer is negligible for these updates.

Start with the width-640 Kaiming baseline from the assignment, including its
embedding, normalization, and readout initialization. Its LR is $.001$ for
\emph{every} parameter group. Instead of independently initializing a wider
model, expand this model by a factor $k$, where $k=4$ or $8$.

Copy the embedding coordinates and width-dependent normalization gains $k$
times. Copy complete attention heads, keeping each head's dimension at 64;
leave the shared 64-dimensional QK-normalization gains unchanged.
For each attention or feedforward matrix $W$ (including the rectangular
gate, up, and down projections), replace it by the
$k$-by-$k$ block matrix whose every block is $W/k$. If $V$ is the readout
matrix, concatenate $k$ copies along its input dimension and divide by $k$:
\[
 W'=\frac1k\begin{bmatrix}W&\cdots&W\\\vdots&\ddots&\vdots\\W&\cdots&W\end{bmatrix},
 \qquad V'=\frac1k[V\ \cdots\ V].
\]
There is no additional readout multiplier. The copies are separate trainable
parameters, not tied weights.

\textbf{Which learning rates should most closely preserve the width-640
model's logits on the same inputs after each of the five updates?}
Select one option for each width. The entries below are actual parameter-group
LRs, with no additional width multipliers. ``Other'' means embedding and all
normalization parameters.

\textbf{(a) Width 2560.}
\begin{center}\begin{tabular}{lrrr}
\toprule
Option & Hidden matrices & Readout & Other\\\midrule
\checkbox\ 1 & $.001$ & $.001$ & $.001$\\
\checkbox\ 2 & $.00025$ & $.00025$ & $.00025$\\
\checkbox\ 3 & $.00025$ & $.00025$ & $.001$\\
\checkbox\ 4 & $.00025$ & $.001$ & $.001$\\
\bottomrule\end{tabular}\end{center}
\textbf{(b) Width 5120.}
\begin{center}\begin{tabular}{lrrr}
\toprule
Option & Hidden matrices & Readout & Other\\\midrule
\checkbox\ 1 & $.00025$ & $.00025$ & $.001$\\
\checkbox\ 2 & $.000125$ & $.000125$ & $.001$\\
\checkbox\ 3 & $.000125$ & $.000125$ & $.000125$\\
\checkbox\ 4 & $.000125$ & $.001$ & $.001$\\
\bottomrule\end{tabular}\end{center}

\end{examplequestion}

